%=====================================================================
\chapter{What Artificial Intelligence Is}\label{ch:introduction}
%=====================================================================
\locator{1}{Part I --- Foundations of Intelligent Systems}

\begin{outcomes}
By the end of this chapter you will be able to:
\begin{tight}
  \item \textbf{Define} artificial intelligence, and \textbf{distinguish} the
        four classical definitions along the two axes of thinking versus acting
        and humanly versus rationally.
  \item \textbf{Identify} whether a given piece of software is an AI system, by
        applying a five-point test rather than by reputation.
  \item \textbf{Name} the branches of AI and \textbf{say} which of them this
        course teaches, which it signposts, and which belong to other courses.
  \item \textbf{Describe} the field's history as a sequence of promises,
        disappointments and recoveries, and \textbf{explain} what caused each
        turn.
  \item \textbf{State} the boundary between artificial intelligence, machine
        learning and data science in terms of what each one is given and what
        each one must work out for itself.
\end{tight}
\end{outcomes}

\begin{prereq}
Object Oriented Programming, the published prerequisite for this course: classes,
objects, inheritance, collections and recursion. No mathematics beyond school
algebra is assumed in this chapter. Appendix~A is a reference for every Python
construct the labs use.
\end{prereq}

\begin{keyterms}
artificial intelligence $\bullet$ agent $\bullet$ percept $\bullet$ action
$\bullet$ rationality $\bullet$ the Turing test $\bullet$ symbolic AI $\bullet$
knowledge representation $\bullet$ search $\bullet$ heuristic $\bullet$
combinatorial explosion $\bullet$ weak and strong AI $\bullet$ narrow and general
AI $\bullet$ AI winter $\bullet$ expert system
\end{keyterms}

\section{Two Programs in a Project Office}

A software house runs its delivery work out of a project management office. Two
programs run there every night.

The first recalculates the Gantt chart. It is given the task list, each task's
duration and each task's predecessors, and it computes the earliest start and
latest finish of every task, the critical path, and the float on everything off
it. It is a few hundred lines long, it is completely deterministic, and if it
ever disagrees with a hand calculation the program is wrong. Nobody in the office
calls it intelligent.

The second decides which of tonight's three hundred open tickets to escalate to
the on-call engineer. It is given the tickets, the service level agreements, who
is on shift, what each engineer has already been handed, and a history of which
escalations turned out to have been worth making. There is no formula in the
brief. Two reasonable people, given the same three hundred tickets, would escalate
different ones and each could defend the choice. The program has to decide anyway,
tonight, before midnight.

That difference --- not the difficulty of the arithmetic, not the size of the
codebase, and not the fashionability of the technique --- is what this book is
about. The first program computes an answer that was already implied by its
input. The second must work out what to do when its input does not determine the
answer.

\begin{definitionbox}[Artificial intelligence]
\term{Artificial intelligence} is the study and construction of \term{agents}:
programs that perceive an environment through \term{percepts} and act on it
through \term{actions}, choosing those actions so as to do well by some measure of
success. The central question of the field is not ``can a machine think?'' but
``given what it can perceive, what should this agent do next?''

\smallskip
\noindent\emph{After} S.~Russell and P.~Norvig, \emph{Artificial Intelligence: A
Modern Approach}, 3rd edition, Prentice Hall, 2010 --- the set text for this
course in both the HEC 2017 and HEC 2023 outlines.
\end{definitionbox}

The definition is deliberately modest. It says nothing about consciousness,
nothing about brains, and nothing about whether the agent's reasoning resembles
ours. It asks only that the agent take in a situation and choose. Everything in
this course is machinery for making that choice: search when the agent must look
ahead (Chapters~4 to~6), constraint reasoning when the choice must satisfy rules
at once (Chapter~7), game reasoning when another agent is choosing against you
(Chapter~8), planning when the choice is a sequence (Chapter~9), logic when the
agent must draw conclusions from what it has been told (Chapters~10 to~13), and
probability when what it has been told is not certain (Chapters~14 to~16).

\section{Four Definitions, Two Axes}

The field has never had one definition. It has had four, and they differ along two
axes: whether the goal is to replicate \emph{thought} or to produce the right
\emph{behaviour}, and whether the standard of correctness is \emph{human}
performance or an ideal standard of \emph{rationality}. \dsref{ai-axes} sets them
out.

\dsfig{%
\begin{tikzpicture}[font=\small,
  qbox/.style={rounded corners=3pt, draw=#1, line width=1pt, fill=#1!7,
               text width=5.6cm, align=flush left, inner sep=7pt,
               minimum height=2.5cm, text=#1!55!black},
  axl/.style={font=\scriptsize\bfseries, text=neutral}]

\node[qbox=purplehl] (tw) at (0,0)
  {{\bfseries Thinking humanly}\\[2pt]
   {\scriptsize The cognitive modelling approach. A program is right if it solves
    the problem by the same steps a person does. Validated against
    psychological experiment, not against the answer.}\\[3pt]
   {\scriptsize\itshape Newell \& Simon, General Problem Solver, 1961}};

\node[qbox=primary, right=0.7cm of tw] (tr)
  {{\bfseries Thinking rationally}\\[2pt]
   {\scriptsize The laws-of-thought approach. Correct reasoning is a formal
    relation between premises and conclusions; build a program that computes it.}\\[3pt]
   {\scriptsize\itshape Aristotle's syllogisms to modern logic ---
    Chapters~10 to~13}};

\node[qbox=goldhl, below=0.7cm of tw] (aw)
  {{\bfseries Acting humanly}\\[2pt]
   {\scriptsize The Turing test approach. A program is intelligent if an
    interrogator cannot tell it from a person in conversation. Behaviour is the
    only evidence admitted.}\\[3pt]
   {\scriptsize\itshape Turing 1950; Eliza, 1966 --- Chapters~18 and~19}};

\node[qbox=greenhl, below=0.7cm of tr] (ar)
  {{\bfseries Acting rationally}\\[2pt]
   {\scriptsize The rational agent approach. Given what it perceives and what it
    is trying to achieve, the agent takes the action that can be expected to do
    best. Human imitation is not required.}\\[3pt]
   {\scriptsize\itshape \bfseries The approach this course takes}};

\node[axl, above=3pt of tw.north] {HUMAN STANDARD};
\node[axl, above=3pt of tr.north] {RATIONAL STANDARD};
\node[axl, rotate=90, left=3pt of tw.west] {REASONING};
\node[axl, rotate=90, left=3pt of aw.west] {BEHAVIOUR};

\draw[greenhl, line width=1.6pt, rounded corners=4pt]
  ($(ar.north west)+(-3pt,3pt)$) rectangle ($(ar.south east)+(3pt,-3pt)$);
\end{tikzpicture}}%
{The four classical definitions of artificial intelligence, on the two axes of
thought against behaviour and human performance against rational performance.
This course works in the bottom-right quadrant throughout.}{ai-axes}

\begin{conceptbox}
The four quadrants are not four opinions about the same question; they are four
different questions, and they license different evidence. If your goal is
\emph{acting humanly}, then a program that answers a service-desk query the way a
tired engineer would --- including the shortcuts --- is a success. If your goal is
\emph{acting rationally}, the same program is a failure, because the tired
engineer was wrong.

This matters in practice, and the practical consequence is this: \emph{decide
which quadrant you are in before you choose your evaluation}. Chapter~19 shows
Eliza, a program that is a triumph in the bottom-left quadrant and worthless in
the bottom-right, and the difference is not in the code.
\end{conceptbox}

\section{Is It AI? A Checklist}
\label{sec:is-it-ai}

``Artificial intelligence'' is now a procurement word. Vendors apply it to
spreadsheet macros, and sceptics withhold it from systems that plainly reason. The
HEC outline asks you to be able to \emph{identify} AI systems, which means having
a test you can apply to software you did not write. Here is one. Five questions;
score each one and count.

\begin{enumerate}[leftmargin=2.2em, itemsep=3pt, topsep=4pt]
  \item \textbf{Goal.} Is the program given an \emph{objective} rather than a
        procedure? A program told \emph{minimise the number of breached service
        level agreements} has a goal. A program told \emph{sort the tickets by
        deadline, then assign the first ten} has a procedure.
  \item \textbf{Search.} Does it \emph{consider alternatives it does not act on}?
        An agent that generates candidate assignments, evaluates them and picks
        one has searched. An agent that computes one answer directly has not.
  \item \textbf{Knowledge.} Does it hold an explicit, inspectable
        \emph{representation} of its domain --- facts, rules, constraints, a
        probability model --- that a domain expert could read and correct without
        editing the algorithm?
  \item \textbf{Adaptation.} Does its behaviour change with \emph{experience},
        without being reprogrammed?
  \item \textbf{Explanation.} Can it say \emph{why}? Can it produce the chain of
        facts and inferences that led to its answer?
\end{enumerate}

A program scoring zero is ordinary software, however valuable. A program scoring
one or two is doing something interesting. Three or more, and the vocabulary of
this course is the right vocabulary for talking about it.

Two warnings about the test. First, it is not a quality measure: a badly built
system can score five, and the Gantt recalculator that scores zero is the most
useful program in the office. Second, question~4 is neither necessary nor
sufficient. Much of the most successful AI ever built --- the chess programs of
Chapter~8, the schedulers of Chapter~7, the diagnostic systems of Chapter~13 ---
scores zero on adaptation. Learning is one route to intelligent behaviour, not its
definition. That is worth saying plainly at the start of a course that will spend
fifteen weeks on the other routes.

\begin{examplebox}[Five tools in one project office]
Scores on the five tests of Section~\ref{sec:is-it-ai}, for the five programs the
office actually runs. This is the audit the chapter's lab reproduces in code.

\smallskip
\noindent\begin{tabular}{@{}L{5.3cm}ccccc@{\hspace{6pt}}c@{}}
\toprule
\textbf{Program} & \textbf{Goal} & \textbf{Srch} & \textbf{Know} &
\textbf{Adapt} & \textbf{Expl} & \textbf{Score} \\
\midrule
Gantt recalculator        & --- & --- & --- & --- & --- & 0 \\
Burndown chart generator  & --- & --- & --- & --- & --- & 0 \\
Sprint scheduler          & yes & yes & yes & --- & yes & 4 \\
Escalation adviser        & yes & yes & yes & --- & yes & 4 \\
Duplicate-ticket finder   & --- & --- & --- & yes & --- & 1 \\
\bottomrule
\end{tabular}

\smallskip
The two zeroes are not failures; they are the right tools for questions that have
formulas. The sprint scheduler is Chapter~7, the escalation adviser is
Chapter~13, and the duplicate finder is the one tool of the five that this course
hands to \chref{learning} and then to Machine Learning.
\end{examplebox}

\section{The Branches of AI}

The field is conventionally divided as in \dsref{ai-branches}. The division is
untidy, because the branches are not disjoint --- planning uses search, natural
language processing uses both logic and probability --- but it is the division the
HEC outline uses and the one job advertisements use.

\dsfig{%
\begin{tikzpicture}[font=\scriptsize,
  root/.style={rounded corners=4pt, fill=primary, text=white,
               font=\small\bfseries, minimum width=3.4cm, minimum height=0.85cm,
               align=center},
  br/.style={rounded corners=3pt, draw=#1, line width=1pt, fill=#1!10,
             text=#1!55!black, font=\scriptsize\bfseries, text width=2.45cm,
             align=center, minimum height=0.8cm},
  sub/.style={font=\tiny, text=neutral, text width=2.45cm, align=center},
  a/.style={-{Stealth[length=1.8mm]}, primary!45, line width=0.8pt}]

\node[root] (r) at (0,0) {Artificial Intelligence};

\node[br=primary]   (b1) at (-6.3,-1.9) {Search \& Problem Solving};
\node[br=secondary] (b2) at (-3.15,-1.9) {Knowledge \& Reasoning};
\node[br=greenhl]   (b3) at (0,-1.9)     {Planning \& Scheduling};
\node[br=tealhl]    (b4) at (3.15,-1.9)  {Uncertainty \& Probability};
\node[br=purplehl]  (b5) at (6.3,-1.9)   {Machine Learning};

\node[sub, below=2pt of b1] {Ch.\ 4--6, 8\\[-1pt]states, heuristics, games};
\node[sub, below=2pt of b2] {Ch.\ 10--13\\[-1pt]logic, rules, ontologies};
\node[sub, below=2pt of b3] {Ch.\ 7, 9\\[-1pt]constraints, STRIPS};
\node[sub, below=2pt of b4] {Ch.\ 14--16\\[-1pt]Bayes nets, fuzzy logic};
\node[sub, below=2pt of b5] {Ch.\ 17 only\\[-1pt]\emph{Machine Learning
      Fall 2026}};

\node[br=goldhl]  (b6) at (-4.7,-4.6) {Natural Language Processing};
\node[br=goldhl]  (b7) at (-1.55,-4.6) {Computer Vision};
\node[br=goldhl]  (b8) at (1.6,-4.6)   {Robotics};
\node[br=goldhl]  (b9) at (4.75,-4.6)  {Multi-Agent Systems};

\node[sub, below=2pt of b6] {Ch.\ 18 (symbolic)};
\node[sub, below=2pt of b7] {\emph{Computer Vision}};
\node[sub, below=2pt of b8] {not in this degree};
\node[sub, below=2pt of b9] {Ch.\ 20 (signpost)};

\foreach \b in {b1,b2,b3,b4,b5}{\draw[a] (r) -- (\b.north);}
\foreach \b in {b6,b7,b8,b9}{\draw[a, goldhl!55] (r.south) -- (\b.north);}

\node[font=\tiny\itshape, text=neutral, align=left, anchor=west]
  at (-6.6,-6.0) {Bold branches are taught here. Grey captions name the course
  that owns the branch instead; Appendix~F states every boundary in full.};
\end{tikzpicture}}%
{The branches of artificial intelligence, and which of them this course teaches.
The five branches on the upper row are the body of the course; the lower row is
covered by a signpost chapter or by another course entirely.}{ai-branches}

\section{Sixty Years in Five Waves}

AI has a distinctive history: repeated cycles in which a genuine advance is
over-promised, fails to generalise, and is followed by a collapse in funding
severe enough to have earned its own name, an \term{AI winter}. Knowing the shape
of the cycle is professionally useful, because it is happening again now and you
will be asked to make judgements inside it.

\smallskip
\noindent\begin{longtable}{@{}L{2.5cm}L{5.4cm}L{7.1cm}@{}}
\toprule
\textbf{Period} & \textbf{What was achieved} & \textbf{Why it stalled} \\
\midrule
\endfirsthead
\toprule
\textbf{Period} & \textbf{What was achieved} & \textbf{Why it stalled} \\
\midrule
\endhead
1943--1956 \newline \emph{Foundations} &
McCulloch and Pitts' artificial neuron; Turing's 1950 paper and his test; the
1956 Dartmouth workshop, which named the field. &
Nothing stalled yet. The founding claim --- that every aspect of intelligence
could be ``so precisely described that a machine can be made to simulate it'' ---
set the expectations that the next twenty years failed to meet. \\
1956--1974 \newline \emph{Search and the general solver} &
The Logic Theorist and the General Problem Solver (Chapter~9); Samuel's checkers
player; Eliza (Chapter~18); early perceptrons (Chapter~17). &
\term{Combinatorial explosion}. The methods worked on toy problems and died on
real ones: a search that is fine at depth 6 is hopeless at depth 20. Minsky and
Papert's 1969 analysis of the perceptron's limits closed the neural line for
fifteen years. \\
1974--1980 \newline \emph{First winter} &
--- &
The Lighthill report (1973) in the UK and DARPA's withdrawal in the US ended
general AI funding. The diagnosis was correct: the field had solved the easy half
of its problems and called it progress. \\
1980--1987 \newline \emph{Expert systems} &
Knowledge, not cleverness, as the engine: MYCIN, XCON, and a commercial industry
in rule-based systems (Chapter~13). Japan's Fifth Generation project. &
The \term{knowledge acquisition bottleneck}. Rules had to be elicited from experts
one at a time, could not be learned, grew brittle at a few thousand rules, and
degraded catastrophically rather than gracefully at the edge of their competence.
\\
1987--1993 \newline \emph{Second winter} &
--- &
The specialised hardware market collapsed; ``AI'' became a word firms avoided in
prospectuses. The research continued under other names --- ``knowledge-based
systems'', ``decision support''. \\
1993--2011 \newline \emph{The quiet, successful period} &
The rational-agent reframing that this course follows; probability and Bayesian
networks (Chapter~15); Deep Blue (1997); statistical machine translation; search
engines and recommenders. AI methods everywhere, under other names. &
It did not stall. It stopped claiming to be AI, which is why this era is the least
famous and the most useful. \\
2012--present \newline \emph{Learning at scale} &
Deep networks, then transformers and large language models; systems that write
fluent text and usable code (Chapter~20). &
Still open. The symbolic questions this course teaches --- can it prove it, can it
explain it, can it guarantee the constraint holds --- are exactly the ones the
current wave answers least well (Chapter~21). \\
\bottomrule
\caption{Five waves and two winters. Each collapse followed a real achievement
that was mistaken for a general one.}\label{tab:ai-history}
\end{longtable}

\begin{alertbox}[The pattern worth taking from the history]
Every winter arrived the same way: a method that worked at small scale was sold as
a method that worked, and the gap was discovered by customers rather than by
researchers. The methods themselves were not discredited --- minimax, resolution
and rule engines are all in production today, in this book's Chapters~8, 12
and~13. What was discredited each time was a claim about their \emph{reach}.

You will be asked to make that kind of claim in your own career, probably about
something you built over a weekend that worked beautifully on twenty examples.
Chapter~21 gives you the vocabulary to make it honestly.
\end{alertbox}

\section{What This Course Teaches, and What It Does Not}

Artificial intelligence, machine learning and data science are three courses in
this degree, taught in that order, and students routinely confuse them because all
three are described in job advertisements with the same six words.
\dsref{ai-scope} draws the boundary in a way that survives contact with a real
question, and Appendix~F states it chapter by chapter.

\dsfig{%
\begin{tikzpicture}[font=\scriptsize,
  crs/.style={rounded corners=4pt, draw=#1, line width=1.4pt, fill=#1!8,
              text width=4.15cm, align=flush left, inner sep=7pt,
              text=#1!50!black},
  hd/.style={rounded corners=2pt, fill=#1, text=white,
             font=\scriptsize\bfseries, inner sep=3pt},
  a/.style={-{Stealth[length=2.4mm]}, neutral, line width=1.1pt}]

\node[crs=primary] (ai) at (0,0)
  {\\[1pt]{\scriptsize\bfseries Given:} the rules of the world --- the tasks and
   their dependencies, the constraints, the operators, the facts.\\[3pt]
   {\scriptsize\bfseries Must work out:} \emph{what to do}. Which sequence of
   actions reaches the goal.\\[3pt]
   {\scriptsize\bfseries Method:} search and inference over an explicit
   representation.\\[3pt]
   {\scriptsize\bfseries Answer comes with:} a proof, a plan, or a trace.};

\node[crs=greenhl, right=0.75cm of ai] (ml)
  {\\[1pt]{\scriptsize\bfseries Given:} examples of inputs paired with the right
   outputs. The rules are \emph{not} given.\\[3pt]
   {\scriptsize\bfseries Must work out:} \emph{the rule itself}, and whether it
   will hold on data never seen.\\[3pt]
   {\scriptsize\bfseries Method:} fitting a model, then measuring how it
   generalises.\\[3pt]
   {\scriptsize\bfseries Answer comes with:} an error estimate.};

\node[crs=goldhl, right=0.75cm of ml] (ds)
  {\\[1pt]{\scriptsize\bfseries Given:} a question somebody actually asked, and
   messy data that half answers it.\\[3pt]
   {\scriptsize\bfseries Must work out:} \emph{what is true}, and what should be
   done about it.\\[3pt]
   {\scriptsize\bfseries Method:} the whole pipeline --- acquisition, cleaning,
   inference, models, communication.\\[3pt]
   {\scriptsize\bfseries Answer comes with:} a decision and a caveat.};

\node[hd=primary] at (ai.north) {ARTIFICIAL INTELLIGENCE \textnormal{$\cdot$ Sem 3}};
\node[hd=greenhl] at (ml.north) {MACHINE LEARNING \textnormal{$\cdot$ Sem 6}};
\node[hd=goldhl]  at (ds.north) {DATA SCIENCE \textnormal{$\cdot$ Sem 7}};

\draw[a] (ai.east) -- node[above, font=\tiny\itshape, text=neutral]
  {prerequisite} (ml.west);
\draw[a] (ml.east) -- node[above, font=\tiny\itshape, text=neutral]
  {prerequisite} (ds.west);

\node[font=\tiny\itshape, text=neutral, text width=13.5cm, align=flush left]
  at (4.5,-3.35)
  {One question, three courses. \emph{Which tickets should be escalated
   tonight?} --- AI writes the rules down as constraints and infers the
   escalation set, with the chain of reasoning attached (Chapter~13). Machine
   Learning is given ten thousand past escalations and their outcomes and learns
   which features predict a worthwhile one. Data Science asks whether
   ``worthwhile'' was measured correctly in the first place, and what the
   escalation policy is doing to the on-call engineers.};
\end{tikzpicture}}%
{Artificial intelligence, machine learning and data science, separated by what
each is given and what each must work out. The distinction is not the technique;
it is where the rules come from.}{ai-scope}

So: this course is \emph{symbolic}. It teaches the agent that is told the rules
and must reason with them. That is not a historical curiosity, for three reasons.
It is where guarantees live --- a constraint solver can prove no engineer is
double-booked, and no learned model can (Chapter~7). It is where explanations live
--- a rule engine can print the chain that produced its answer, and that chain is
what an auditor accepts (Chapter~13). And it is where you go when you have no
data, which is the situation at the start of every new project.

Chapter~17 is the one chapter about learning. It covers the three paradigms the
HEC 2017 outline names, works the perceptron numerically, shows why a single
perceptron cannot represent exclusive-or --- the result that caused the first
neural winter --- and then stops and hands over. It is a signpost, not a unit.

\begin{worked}[Auditing one tool against the five tests]
The office's \emph{duplicate-ticket finder} takes each new ticket and flags the
three existing tickets most similar to it, so the service desk can merge them. Its
similarity measure is tuned monthly against the merges the desk actually made.
Score it.

\smallskip
\textbf{1. Goal?} \emph{No.} The brief is a procedure: compute similarity, return
the top three. There is no objective it is trying to do well by; if the desk
wanted the top five it would be a different program.
\smallskip

\textbf{2. Search?} \emph{No.} It scores every candidate and returns the best
three. Scoring all alternatives and ranking them is not search --- there is no
point at which the program commits to a partial choice and explores its
consequences. Compare Chapter~4, where a decision about task~3 changes which
states are reachable at all.
\smallskip

\textbf{3. Knowledge?} \emph{No.} The similarity weights are numbers in a
configuration file. A service-desk manager cannot read them and say ``that is
wrong, a database timeout is not like a login failure''.
\smallskip

\textbf{4. Adaptation?} \emph{Yes.} Behaviour changes monthly on the evidence of
past merges, with nobody editing the code.
\smallskip

\textbf{5. Explanation?} \emph{No.} It reports a similarity score. A score is a
number, not a reason.
\smallskip

\textbf{Score: 1 of 5} --- and the one point is adaptation, the single test that
this course does not teach. That is the correct verdict, and the useful one: the
duplicate finder is a good tool and it is not an AI problem in the sense of this
course. It is the machine learning course's problem. Contrast the sprint
scheduler, which scores 4 and whose one missing point is adaptation.
\end{worked}

\begin{pitfall}
\begin{tight}
  \item \textbf{Treating ``AI'' as a synonym for ``machine learning''.} It is the
        most common error in this subject, and it makes whole categories of
        problem invisible. If a constraint must \emph{never} be violated, no
        amount of training data is the answer; you need Chapter~7.
  \item \textbf{Reading the four quadrants of \dsref{ai-axes} as a ranking.} They
        are different goals with different evidence. A program can be excellent
        in one and worthless in the neighbouring one.
  \item \textbf{Assuming a program that scores well on the checklist is a good
        program.} The checklist identifies the \emph{kind} of system, not its
        quality. The highest-scoring system in an office can be the one that
        should be deleted.
  \item \textbf{Believing the winters were caused by the methods failing.}
        Minimax, resolution and rule engines all still work and are all still
        deployed. What failed was a claim about how far they reached.
  \item \textbf{Calling a system ``general'' AI because it is fluent.} Fluency is
        the bottom-left quadrant --- acting humanly. Chapter~19's Eliza is
        two hundred lines long and was mistaken for a therapist in 1966.
\end{tight}
\end{pitfall}

%---------------------------------------------------------------------
\begin{lab}[Is It AI? A Five-Point Audit]
The audit of the \emph{examplebox} above, in code. Five tools from the project
management office, five tests, one verdict each. Pure standard library: the point
of this lab is the classification, not the computation.
\begin{lstlisting}[language=Python]
"""Chapter 1 lab -- Is it AI? A five-point audit.

Score five programs from a project management office against the five
tests of Section 1.3, and classify each one. Nothing here learns
anything; the exercise is to separate programs that decide from
programs that compute.
"""

TESTS = ("goal", "search", "knowledge", "adapts", "explains")

# One row per tool: which of the five tests it passes, and why the
# most interesting answer is what it is.
TOOLS = {
    "Gantt recalculator": {
        "passes": (),
        "note": "given a procedure, not an objective; one answer, "
                "computed directly",
    },
    "Burndown chart generator": {
        "passes": (),
        "note": "arithmetic and a plot; no alternatives considered",
    },
    "Sprint scheduler": {
        "passes": ("goal", "search", "knowledge", "explains"),
        "note": "constraints are declared, not coded; reports which "
                "constraint blocked a story  (Chapter 7)",
    },
    "Escalation adviser": {
        "passes": ("goal", "search", "knowledge", "explains"),
        "note": "rules an operations lead can read and correct; prints "
                "its firing chain  (Chapter 13)",
    },
    "Duplicate-ticket finder": {
        "passes": ("adapts",),
        "note": "weights retuned monthly from past merges; the one "
                "tool here that belongs to Machine Learning",
    },
}


def verdict(score):
    """The reading of a score given in Section 1.3."""
    if score == 0:
        return "ordinary software"
    if score <= 2:
        return "doing something interesting"
    return "an AI system"


def audit():
    width = max(len(name) for name in TOOLS)
    head = "TOOL".ljust(width) + "  " + " ".join(
        t[:4].upper().rjust(5) for t in TESTS) + "   SCORE  VERDICT"
    print(head)
    print("-" * len(head))
    for name, row in TOOLS.items():
        marks = ["yes".rjust(5) if t in row["passes"] else "-".rjust(5)
                 for t in TESTS]
        score = len(row["passes"])
        print(name.ljust(width) + "  " + " ".join(marks) +
              "   %5d  %s" % (score, verdict(score)))
    print()
    for name, row in TOOLS.items():
        print("%s\n    %s" % (name, row["note"]))


def check():
    """The two claims the chapter makes about this table."""
    scores = {n: len(r["passes"]) for n, r in TOOLS.items()}
    # Adaptation is the only test the AI systems here fail ...
    for name in ("Sprint scheduler", "Escalation adviser"):
        assert "adapts" not in TOOLS[name]["passes"], name
        assert scores[name] == 4, name
    # ... and the only test the machine learning tool passes.
    assert TOOLS["Duplicate-ticket finder"]["passes"] == ("adapts",)
    print("\nchecks passed: intelligent behaviour here does not "
          "require learning,")
    print("and the one tool that learns is the one this course "
          "hands on.")


if __name__ == "__main__":
    audit()
    check()
\end{lstlisting}
Now add a sixth tool of your own --- something you have actually used, a CI
pipeline or an IDE's autocomplete --- and score it. The interesting cases are the
ones that score 2 or 3: write down which test you were least sure about and why.
\end{lab}

%---------------------------------------------------------------------
\begin{checkpoint}
\begin{tightnum}
  \item A program is given the sentence \emph{minimise total project delay} and a
        list of tasks with durations and dependencies. Which of the five tests of
        Section~\ref{sec:is-it-ai} does it pass on that description alone, and
        which cannot be judged until you see the implementation?
  \item A chatbot answers service-desk questions so convincingly that users
        believe it is a person, but it has no representation of the service catalogue
        and cannot say why it gave an answer. Which quadrant of \dsref{ai-axes} is
        it succeeding in, and which is it failing in?
  \item The expert-systems boom ended in the knowledge acquisition bottleneck.
        Name the chapter of this handout that builds an expert system, and state
        which single one of the five tests such a system fails.
\end{tightnum}
\end{checkpoint}

%---------------------------------------------------------------------
\begin{chaptersummary}
\begin{tight}
  \item An \term{agent} perceives an environment and acts on it. Artificial
        intelligence is the study of choosing those actions well; the field's
        question is ``what should this agent do next?''.
  \item The four classical definitions differ on two axes: thinking against
        acting, and the human standard against the rational standard. This course
        works throughout in the \emph{acting rationally} quadrant.
  \item A program is an AI system to the extent that it is given a goal rather
        than a procedure, considers alternatives, holds an inspectable
        representation of its domain, adapts, and can explain itself. Three or
        more of the five, and this course's vocabulary applies.
  \item Learning is one route to intelligent behaviour, not the definition of it.
        Chess programs, schedulers and diagnostic systems score zero on
        adaptation and are unambiguously AI.
  \item The history is a cycle: a real advance, an over-general claim, a winter.
        Combinatorial explosion ended the first wave; the knowledge acquisition
        bottleneck ended the expert-system wave.
  \item Artificial intelligence is given the rules and must work out what to do;
        machine learning is given examples and must work out the rule; data
        science is given a question and must work out what is true. This course is
        the first of the three, and it is symbolic.
\end{tight}
\end{chaptersummary}

\begin{reviewq}
  \item \textbf{(MCQ)} A program that beats a grandmaster at chess but cannot
        learn from its games is: (a)~not AI, because it does not learn (b)~AI, and
        scores zero on the adaptation test (c)~AI only in the ``thinking humanly''
        sense (d)~machine learning.
  \item \textbf{(MCQ)} The Turing test belongs in which quadrant?
        (a)~thinking humanly (b)~thinking rationally (c)~acting humanly
        (d)~acting rationally.
  \item \textbf{(MCQ)} Which ended the first wave of AI research?
        (a)~the knowledge acquisition bottleneck (b)~combinatorial explosion
        (c)~the collapse of the Lisp machine market (d)~Minsky and Papert's book
        alone.
  \item \textbf{(MCQ)} A system is given the tasks, the dependencies and the
        objective, and returns a sequence of actions with a proof that no
        dependency is violated. Which course owns this problem?
        (a)~Artificial Intelligence (b)~Machine Learning (c)~Data Science
        (d)~Data Mining.
  \item \textbf{(Short)} State the difference between being given a goal and being
        given a procedure, using two versions of the same ticket-triage
        requirement as your example.
  \item \textbf{(Short)} Why is ``it adapts'' neither necessary nor sufficient for
        a program to be an AI system? Give one example each way.
  \item \textbf{(Short)} The expert-system industry collapsed even though rule
        engines work. Explain what actually failed, and name the chapter of this
        handout in which you will build one anyway.
  \item \textbf{(Applied)} Take a piece of software you have used this week. Score
        it against the five tests, writing one sentence of justification per test,
        and give the verdict. Then state which single change to the system would
        raise its score by one --- and whether that change would make it better.
\end{reviewq}
